\section{相关工作、研究缺口与本文定位}
\subsection{点接触与互补类硬接触方法}
在多数多体动力学建模框架中，接触仍主要依赖点接触或少量离散接触点实现。该类模型在球—平面、简单凸体碰撞等场景中具有较高效率，但在以下问题中存在明显不足：
\begin{enumerate}[label=(\arabic*)]
    \item 难以描述面接触和接触斑演化；
    \item 难以准确给出净接触力矩；
    \item 对接触点采样和几何离散较敏感；
    \item 在复杂几何和多接触切换中易出现抖动。
\end{enumerate}

因此，点接触模型虽然计算简单，但难以满足复杂刚体接触建模对几何表达和力学合理性的双重要求。

互补类方法在理想刚体接触问题中具有重要理论地位，其核心思想是通过非穿透约束和接触力条件构造硬约束模型。然而，互补类方法通常建立在点接触、刚体假设及理想摩擦模型基础上。当接触区域较大、材料存在局部顺应性、表面几何复杂或需要显式求取分布压力与接触力矩时，这类方法难以提供足够细致的接触描述。此外，复杂摩擦、多接触切换以及接触模式不确定性，也会在理论和数值上引入无解、多解或高敏感性问题。因此，互补类方法更适合描述宏观刚体约束，而非任意复杂接触场景下的真实空间分布作用。

\subsection{顺应接触模型与复杂几何接触}
顺应接触模型通过将接触视为局部压缩过程，以穿透量或压入量构造法向力和阻尼力，因此在数值上更连续、实现上更灵活，也更适合工程仿真和复杂几何。然而，传统顺应接触模型大多停留在点级罚函数框架内，未能从根本上解决面接触和分布力矩建模问题。换言之，顺应接触虽然在力学处理方式上比硬约束更柔性，但在空间表达上仍然较为粗糙。

\subsection{SDF 在复杂接触建模中的潜力}
SDF 以距离场形式统一表示几何边界，对任意复杂、非凸几何都具有良好的适应性。其优势包括：
\begin{enumerate}[label=(\arabic*)]
    \item 能直接提供局部距离与法向；
    \item 能定义压入量和近接触区域；
    \item 便于提取局部曲率、法向质量等几何特征；
    \item 适合与隐式建模、可微优化、多分辨率表示结合。
\end{enumerate}

然而，目前若仅将 SDF 用作“更方便的碰撞检测工具”，仍然无法充分解决复杂接触中的分布式作用问题。真正值得研究的方向，是让 SDF 不只是给出“一个接触点”，而是进一步定义“一个接触斑”和“一个压力场”。

\subsection{研究缺口与本文定位}
综上，现有方法的核心缺口不在于“是否能够检测到接触”，而在于“如何在多体动力学框架下，以可计算的代价稳定表达复杂几何之间的接触区域、压力分布及其合力矩”。基于这一认识，本文聚焦如下问题：在不引入完整连续体接触求解的前提下，能否构建一种基于符号距离场的接触斑—压力场统一模型，使其在任意复杂刚体几何之间在线生成分布式接触力与力矩，并兼顾几何表达能力、力学合理性和工程可实现性？

本文的定位不是将 SDF 作为单纯的几何查询工具，而是将其作为接触斑构造、局部压入描述和分布式力学闭合的统一基础。由此，本文试图在传统点接触模型和高代价连续体接触模型之间建立一条适用于机构、机器人和一般机械系统动力学分析的中尺度建模路径。
